To investigate the impact of the pandemic on professional athletics, an exploratory data analysis was performed to identify structural and qualitative shifts in performance. To gain an initial understanding of the dataset, a heatmap-based plot was employed. This visualization maps the aggregate frequency of competitive entries across a grid of months and years, offering a clear view of seasonal activity levels and data density.

To isolate the specific impact of the pandemic year, the dataset was partitioned into two distinct cohorts. The baseline group consists of all pooled IAAF scores recorded between 2015 and 2019 and the other group consists of all IAAF scores recorded in 2020. This aggregation serves as a stable reference for "normal" pre-pandemic performance levels. This analysis was performed to investigate a potential general trend across the investigated groups. 
To evaluate the directionality and magnitude of the performance shift, the relative percentage difference of the median was calculated between the two time periods. Through this approach, the exact degree of performance deviation relative to the pre-pandemic baseline was precisely quantified.

To evaluate the impact of the pandemic on general performance density, we analyzed the dispersion of IAAF scores. This allows us to determine whether the competitive field became more spread out or remained stable during the 2020 season. The interquartile range (IQR) was chosen as the main measure of dispersion instead of the standard deviation (STD) because it is more robust when analyzing elite athletic data. High-performance scores often include "outliers," extraordinary performances that can greatly distort the standard deviation.
The primary metric for measuring the performance spread is the relative percentage change of IQR in 2020 compared to the pooled baseline from 2015 to 2019.